\documentclass[10pt]{article}

\usepackage[english]{babel} 
\usepackage[T1]{fontenc}
\usepackage{cite, url,color} % Citation numbers being automatically sorted and properly "compressed/ranged". 
%\usepackage{pgfplots}
\usepackage{graphics,amsfonts}
\usepackage[pdftex]{graphicx}
\usepackage[cmex10]{amsmath}
% Also, note that the amsmath package sets \interdisplaylinepenalty to 10000
% thus preventing page breaks from occurring within multiline equations. Use:
 \interdisplaylinepenalty=2500
% after loading amsmath to restore such page breaks as IEEEtran.cls normally does.

\usepackage[utf8]{inputenc}
%% Useful packages for creation of two-column and more complex figures

% Compact lists
\usepackage{enumitem}


\usepackage{array}
% http://www.ctan.org/tex-archive/macros/latex/required/tools/
\usepackage{mdwmath}
\usepackage{mdwtab}
%mdwtab.sty	-- A complete ground-up rewrite of LaTeX's `tabular' and  `array' environments.  Has lots of advantages over
%		   the standard version, and over the version in `array.sty'.
% *** SUBFIGURE PACKAGES ***
\usepackage[tight,footnotesize]{subfigure}

\usepackage[top=2cm, bottom=2cm, right=1.6cm,left=1.6cm]{geometry}
\usepackage{indentfirst}

\usepackage{times}
%\usepackage[active]{srcltx}
\graphicspath{{./figure/}}

\setlength\parindent{0pt}
\linespread{1}

\def\C#1{\mathcal{#1}}

\renewcommand{\phi}{\varphi}
% COMMON OPERATORS
\newcommand{\floor}[1]{{\left\lfloor #1\right\rfloor}}
\newcommand{\ceil}[1]{{\left\lceil #1\right\rceil}}
\newcommand{\E}[1]{\mathop{\rm E}\nolimits\left[#1\right]} % Expectation
\newcommand{\pr}[1]{\Pr\left[ #1 \right]}
\renewcommand{\P}[1]{\mathrm{P}\left[ #1 \right]}
\newcommand{\mbb}{\mathbb}
\newcommand{\mrr}{\mathrm}
\newcommand{\mmin}[1]{{\min\left\{#1\right\}}}
\newcommand{\mmax}[1]{{\max\left\{#1\right\}}}
\newcommand{\id}[1]{\mathbf{1}\(#1\)} % Unit function  
\newcommand{\Heav}[1]{\mathcal{H}\left(#1\right)} % Heaviside function
\newcommand{\eps}{\varepsilon}
\newcommand{\rect}[1]{\mathrm{rect}\!\left(#1\right)}
\newcommand{\sinc}[1]{\mathrm{sinc}\!\left(#1\right)}
\newcommand{\bin}[2]{{\left(\begin{array}{c}#1\\#2\end{array}\right)}}

\newcommand{\me}[2]{\left[ {#1} \right]_{(#2)}} % Submatrix \me{P}{i,j} produces [P]_{i,j} and denotes the element in the i-th row and j-th column of P
\newcommand{\vv}[1]{\left[ {#1} \right]} % Submatrix \vv{a,b,c} produces [a,b,c]

\newcommand{\Set}[1]{{\C #1}} 
\newcommand{\Setd}[2]{\Set{#1}=\left\{#2\right\}}  % Set definition: \Set{S}{0,1,2} produces S={01,2,} where S is in mathcal 

\newcommand{\mat}[1]{{\hbox{\textbf{#1}}}}
\newcommand{\ei}[1]{{\mat{e}_{#1}}}   % all zero vector except in the $#1$-th element which is one
\newcommand{\ind}[1]{\mathbf{\chi}\left\{#1\right\}} % Indicator function \ind{A}=1 if A is true, \ind{A}=0 otherwise. 

% FORMATTING
\newcommand{\ie}{i.e.,\,}
\newcommand{\eg}{e.g.,\,}
\newcommand{\columnbreak}{\vfill\eject} % Column break

% REFERENCES
\newcommand{\Fig}[1]{Fig.~\ref{#1}}
\newcommand{\eq}[1]{(\ref{#1})}
\newcommand{\Tab}[1]{Tab.~\ref{#1}}
\newcommand{\Sec}[1]{Sec.~\ref{#1}}

% COMMENTS/REVISIONS
\newcommand{\AZ}[1]{\textit{\color{magenta}#1}} 

\begin{document}
\title{Communication Network Design  - AY 2020/2021\\
Homework \AZ{replace this \\AZ command with Homework number and title}}
\author{\AZ{replace this with your name and student ID}}
\date{}
\maketitle
\section{Development}
\AZ{Write here your report}.

Example of not-numbered equation
$$
\P{x\leq a}=1-e^{-\mu a}\;.
$$

Example of  numbered equation \eqref{eq:exponentialCDF}
\begin{equation}
\P{x\leq a}=1-e^{-\mu a}\;.
\label{eq:exponentialCDF}
\end{equation}

Example of array of equations \eqref{eq:2rows}
\begin{align}
\P{x\leq a}&=1-e^{-\mu a}\;;  \cr
\P{x\leq a}&=\mmin{a,T}/T\;. 
\label{eq:2rows}
\end{align}


In Fig.~\ref{fig:examplefig} we report an example of figure: 
\begin{figure}[htbp]
\begin{center}
\includegraphics[width=0.5\textwidth]{TestFig.eps}
\caption{Example of figure}
\label{fig:examplefig}
\end{center}
\end{figure}

\AZ{Using this command for the bibliography, you need to check the order of the entries manually, which is fine for short documents. Should you use LaTex for longer documents,  I recommend the use of another bibliography environment called BibTex. See \url{https://en.wikibooks.org/wiki/LaTeX/Bibliography_Management} for more details.}
\begin{thebibliography}{9}

\bibitem{latexcompanion} 
Michel Goossens, Frank Mittelbach, and Alexander Samarin. 
\textit{The \LaTeX\ Companion}. 
Addison-Wesley, Reading, Massachusetts, 1993.

\end{thebibliography}\end{document}
